\documentclass[11pt]{article}
\usepackage[a4paper,margin=24mm]{geometry}
\usepackage{amsmath,amssymb,amsthm,booktabs}
\usepackage[hidelinks]{hyperref}
\newtheorem{theorem}{Theorem}
\newtheorem{proposition}[theorem]{Proposition}
\newtheorem{definition}[theorem]{Definition}
\newcommand{\R}{\mathbb R}
\title{The Q7 game: claims, counterexamples and repair\\
\large For the Q7 resource-allocation game}
\author{}
\date{4 November 2026}
\begin{document}
\maketitle
\vspace{-2em}

Garjani, Shokri and Kebriaei~\cite{q7} propose a resource-allocation
game with quadratic payments. We define the game, restate its claimed
guarantees, and disprove uniqueness and individual rationality by explicit
examples. Removing an opponent-only payment term restores individual
rationality; uniqueness must instead be weakened to uniqueness of the
allocation.

\section{The game and the payment}
There are $N\ge2$ players and $K$ resources with capacity
$c\in\R_+^K$. Player $n$ has a compact convex set
$X_n\subset\R_+^K$ containing zero and a private valuation
$V_n(x_n)=\mathbb E[\psi_n(x_n;\xi_n)]$, with $V_n(0)=0$. We use continuously differentiable, strongly concave
valuations; both counterexamples below are smooth and $1$-strongly concave.
All vector inequalities are coordinatewise, $\|\cdot\|$ is the Euclidean
norm, and $a^\top b$ is the usual inner product.
The planner's problem is
\begin{equation}\label{planner}
 \max_{x_n\in X_n}\sum_{n=1}^N V_n(x_n)
 \quad\text{subject to}\quad \sum_nx_n\le c.
\end{equation}
Player $n$ submits a quantity and a nonnegative price vector,
$s_n=(x_n,p_n)\in X_n\times\R_+^K$. The allocation is exactly $x_n$;
shared capacity is not a constraint on unilateral messages.
Write $s=(s_1,\ldots,s_N)$ for the full message profile and $s_{-n}$
for all messages except player $n$'s.
Utility is $U_n(s)=V_n(x_n)-t_n(s)$: positive tax is paid by the player,
and negative tax is a subsidy. A player submits neither its valuation
function nor a report of its type.

\begin{definition}
A profile $s^*$ is a Nash equilibrium if
$U_n(s_n^*,s_{-n}^*)\ge U_n(s_n,s_{-n}^*)$ for every player and every
admissible joint quantity--price deviation $s_n$.
Individual rationality means $U_n(s^*)\ge V_n(0)=0$ for every player.
Budget balance means $\sum_nt_n(s^*)=0$.
An equilibrium is efficient if its allocation solves~\eqref{planner}.
\end{definition}

Individual rationality here is an equilibrium participation guarantee
relative to an outside option of zero. It does not guarantee that a
learning process reaches equilibrium or that participation is worthwhile
along the way. Randomness enters the static game only through the expected
valuation $V_n$; deterministic valuations are admissible special cases.

Write the paper's payment scale as $\gamma>0$ (its notation is
$\alpha$). Define four separate components, with the first subscript
identifying the player and the second identifying the payment term:
\begin{align}
 t_{n,1}(s)&=\frac{\gamma}{2(N-1)}
 \sum_{m\ne n}\|p_n-p_m\|^2,
 &&\text{price disagreement},\label{term1}\\
 t_{n,2}(s)&=-\frac{\gamma}{N-1}p_n^\top
 \left(\sum_mx_m-c\right),
 &&\text{excess-demand signal},\label{term2}\\
 t_{n,3}(s)&=\frac{\gamma N}{(N-1)^2}
 \left(\sum_{m\ne n}p_m\right)^\top(x_n-c/N),
 &&\text{quantity charge},\label{term3}\\
 t_{n,4}(s)&=-\frac{\gamma}{(N-1)^2}
 \sum_{m\ne n}p_m^\top(x_m-c/N),
 &&\text{opponent-only adjustment}.\label{term4}
\end{align}
The original payment is
\begin{equation}\label{original}
 t_n^{\rm orig}(s)=t_{n,1}(s)+t_{n,2}(s)+t_{n,3}(s)+t_{n,4}(s).
\end{equation}
The first term alone makes the best price the others' average.
The first two together, holding quantities fixed, give
\[
 p_n^{\rm BR}=\left[
 \frac1{N-1}\sum_{m\ne n}p_m+
 \frac{\sum_mx_m-c}{N-1}\right]_+,
\]
where projection onto the nonnegative orthant is coordinatewise.
Excess demand therefore raises proposed prices. The second and third
terms together give
\[
 \nabla_{x_n}t_n^{\rm orig}
 =\gamma\left(-\frac{p_n}{N-1}
 +\frac{N}{(N-1)^2}\sum_{m\ne n}p_m\right)
 =\gamma p\quad\text{when all }p_m=p.
\]
The reference $c/N$ changes transfer levels, not marginal quantity
incentives. The last term affects the recipient's utility but not its
best response, since it depends only on opponents' messages.

\section{The Q7 theorem being claimed}
In the subsection entitled ``A feasible solution to linear matrix
inequalities'', the proposition introducing the explicit payment
(source label \texttt{prop2}) asserts the following, restated in this
note's notation. This is the explicit-mechanism claim we test, rather
than a claim that arbitrary quadratic payments work.
\begin{quote}
\textbf{Claim (Q7, explicit payment).}
The payment $t_n^{\rm orig}$ in~\eqref{original} satisfies the paper's
sufficient matrix conditions for P1--P4, and consequently the induced
mechanism has: P1, an existing and unique Nash equilibrium; P2, an
equilibrium allocation solving the welfare problem; P3, zero total
equilibrium payments; and P4, nonnegative equilibrium utility for every
player relative to the zero outside option.
\end{quote}
These are claims of the original paper, not assumptions of the
counterexample below. Stochastic learning convergence is the separate
property P5, treated subsequently in the paper.

\section{Why the claim is false}
\subsection{Failure of uniqueness: a continuum of price equilibria}
\begin{proposition}\label{nonunique}
The original game can have infinitely many Nash equilibria, even with
a unique efficient allocation and increasing, strongly concave valuations.
\end{proposition}
\begin{proof}
Set $N=3$, $K=1$, $c=3$, $\gamma=1$, and, for every player,
\[
 X_n=[0,1],\qquad V_n(x)=2x-\tfrac12x^2.
\]
For every $q\in[0,1]$, consider the profile
\[
 x_1=x_2=x_3=1,\qquad p_1=p_2=p_3=q.
\]
At this profile, differentiation of the original utility gives
\[
 \partial_{x_n}U_n=1-q\ge0,\qquad
 \partial_{p_n}U_n=0.
\]
Each quantity is at its upper boundary. The own utility Hessian,
in $(x_n,p_n)$, is
\[
 \begin{pmatrix}-1&1/2\\1/2&-1\end{pmatrix}\prec0.
\]
Thus the boundary first-order conditions suffice for global best
responses, including at $q=0$. Every displayed profile is a Nash
equilibrium. All four payment terms vanish: prices agree, total
demand equals capacity, and $x_n=c/N=1$. Each player has utility
$V_n(1)=3/2$.
\end{proof}

This example directly disproves the uniqueness conclusion in
\textbf{P1} for the explicit payment. It preserves existence,
efficiency (P2), budget balance (P3), and participation (P4).
All equilibria displayed have the same allocation but different price
messages. Strict concavity in each player's own strategy ensures a
unique best response to fixed opponents, not a unique joint equilibrium.
Deleting the opponent-only term preserves this continuum, so the
participation repair below does not repair full-equilibrium uniqueness.
The corrected theorem below proves uniqueness only of the allocation.


\subsection{Failure of individual rationality: negative equilibrium utility}
\begin{proposition}\label{counterexample}
The explicit original payment need not be individually rational, even
with smooth strongly concave deterministic valuations and an efficient,
budget-balanced equilibrium.
\end{proposition}
\begin{proof}
Set $N=3$, $K=1$, $c=\gamma=1$, $X_n=[0,2]$, and
\[
 V_1(x)=5x-\tfrac12x^2,\qquad
 V_2(x)=V_3(x)=-x-\tfrac12x^2.
\]
These valuations are normalized and $1$-strongly concave. Deterministic
valuations satisfy the stochastic model as a special case. Consider
\[
 x^*=(1,0,0),\qquad p^*=(4,4,4).
\]
For every player the Hessian of its own utility in $(x_n,p_n)$ is
\[
 H=\begin{pmatrix}-1&1/2\\1/2&-1\end{pmatrix},
 \qquad \operatorname{eig}(H)=\{-3/2,-1/2\}.
\]
Thus own utility is strictly concave. At this profile the own gradients
are $(0,0)$ for player 1 and $(-5,0)$ for players 2 and 3.
The latter players are at the lower quantity boundary. These conditions
prove global best responses over $[0,2]\times\R_+$, hence a Nash
equilibrium. The allocation also solves~\eqref{planner}: players 2 and
3 have strictly decreasing valuations, while player 1's valuation is
increasing throughout $[0,1]$.

Direct substitution in~\eqref{original} gives
\[
\begin{array}{c|rrrr|rrr}
 &t_{n,1}&t_{n,2}&t_{n,3}&t_{n,4}&t_n^{\rm orig}&V_n&U_n\\\hline
 1&0&0&4&2/3&14/3&9/2&-1/6\\
 2&0&0&-2&-1/3&-7/3&0&7/3\\
 3&0&0&-2&-1/3&-7/3&0&7/3
\end{array}
\]
Transfers balance, but player 1's utility is negative. It would prefer
the stipulated outside option of zero. This does not contradict its
best response within the game: requesting zero is still subject to the
payment rule, whereas opting out is the external participation benchmark.
\end{proof}

\paragraph{Exactly what this disproves.}
The example directly contradicts \textbf{P4}: player 1 receives $-1/6$
at a Nash equilibrium. Thus the original proposition's combined
P1--P4 conclusion is false. This example does not address the failure of uniqueness; that is
established separately above.
It satisfies \textbf{P2} and \textbf{P3}: allocation
is optimal and transfers sum to zero. It therefore does not disprove
either of those properties. Nor does a static example alone disprove
the separate convergence claim P5. A player can have negative utility
even if a learning process converges perfectly to this equilibrium.

The claimed satisfaction of the sufficient conditions also fails
directly. In the original parameterization, the linear-in-price
part of the tax is $\sum_m p_m^\top a_m^n$. The paper's sufficient
participation conditions include $\sum_m a_m^n\le0$.
Collecting the capacity terms in~\eqref{original} instead gives
\[
 \sum_m a_m^n
 =\frac{\gamma c}{N-1}-\frac{\gamma c}{N}
 =\frac{\gamma c}{N(N-1)}.
\]
This is positive on every positive-capacity coordinate; in the example
it is $1/6$. Hence the proposition's assertion that its explicit
payment satisfies the sufficient participation conditions is false.
This does not refute the separate conditional participation proposition
for coefficients that actually satisfy those conditions.

\section{How to repair the mechanism and the theorem}
The correction has two parts. Change the payment to restore individual
rationality, and replace the claim of a unique message equilibrium by a
claim of a unique equilibrium allocation. The payment change cannot
select a unique price, because it leaves all best responses unchanged.

\subsection{Payment repair}
Define the repaired payment by keeping just the first three terms:
\begin{equation}\label{repair}
 t_n^c(s)=t_{n,1}(s)+t_{n,2}(s)+t_{n,3}(s)
 =t_n^{\rm orig}(s)-t_{n,4}(s).
\end{equation}
All unilateral utility differences are unchanged. Consequently the
original and repaired games have exactly the same best responses and
Nash equilibria; only utility levels and transfers change.

\paragraph{Why was the fourth term included?}
After presenting the payment, Q7 says that the additional components
depending on opponents' messages are needed to satisfy the LMI conditions.
It gives no separate economic justification for this particular adjustment.
Such terms can redistribute utility without changing incentives. Here,
however, the asserted conditions fail, and the first three terms already
balance at equilibrium. The fourth term is therefore unnecessary for
that balance and can destroy individual rationality, as the example shows.

\subsection{Corrected theorem and proof}
\begin{theorem}[Static implementation and participation]\label{repair-theorem}
Assume each $V_n$ is continuously differentiable and $m$-strongly concave
on $X_n$, with $V_n(0)=0$, and $\gamma\le m(N-1)^2$.
Assume the planner has a primal-dual optimum $(x^*,\lambda^*)$
satisfying its KKT conditions. Then the repaired game has an equilibrium,
and every equilibrium has the planner's unique optimal allocation,
common prices, exact budget balance and individual rationality.
Equilibrium prices need not be unique.
\end{theorem}
\begin{proof}
Put $S=\sum_nx_n$. The negative own price gradient is
\[
 F_{p_n}=\nabla_{p_n}t_n^c
 =\gamma\left[p_n-\frac1{N-1}\sum_{j\ne n}p_j
 -\frac{S-c}{N-1}\right].
\]
At any equilibrium, orthant optimality gives
$F_{p_n}\ge0$ and $p_n^\top F_{p_n}=0$.
Summing the gradients yields
$\sum_nF_{p_n}=\gamma N(c-S)/(N-1)$, hence $S\le c$.
For $\bar p=N^{-1}\sum_np_n$ and $r=c-S\ge0$,
summed complementarity reads
\[
 0=\frac{\gamma N}{N-1}\sum_n\|p_n-\bar p\|^2
 +\frac{\gamma}{N-1}\left(\sum_np_n\right)^\top r.
\]
Both terms are nonnegative. Therefore $p_n=p$ for all $n$ and
$p^\top(c-S)=0$. Allocation optimality becomes
\[
 \langle-\nabla V_n(x_n)+\gamma p,z_n-x_n\rangle\ge0
 \qquad(z_n\in X_n).
\]
Together these are precisely the planner KKT conditions with
$\lambda=\gamma p$, so every equilibrium allocation is optimal and
is unique by strong concavity.

Conversely, set every price to $\lambda^*/\gamma$ at a planner
primal-dual optimum. The displayed first-order conditions hold.
They suffice for global best responses because each player's negative
utility has own Hessian bounded below by
\[
 \begin{pmatrix}
 mI&-\gamma I/(N-1)\\
 -\gamma I/(N-1)&\gamma I
 \end{pmatrix}\succeq0,
\]
by $\gamma\le m(N-1)^2$. For continuously differentiable strongly
concave valuations the same conclusion follows by subtracting the
strong-convexity quadratic, without requiring second derivatives.
This establishes existence.

At common prices the repaired transfer simplifies to
\[
 t_n^c=\gamma p^\top
 \left(x_n-\frac1{N-1}\sum_{j\ne n}x_j\right),
\]
so summing gives zero. Player $n$ can deviate to $(0,p)$, whose tax is
$-\gamma p^\top\sum_{j\ne n}x_j/(N-1)\le0$.
Its utility from that deviation is nonnegative. Equilibrium optimality
therefore gives $U_n^c(s^*)\ge0$.
\end{proof}

In the individual-rationality counterexample, repaired taxes are $(4,-2,-2)$ and utilities
are $(1/2,2,2)$. The equilibrium is unchanged, transfers balance, and
all players benefit from participation. In the nonuniqueness example,
the entire family indexed by $q\in[0,1]$ remains an equilibrium family.
Thus P1 is corrected to existence and allocation uniqueness, P4 is
restored by the payment change, and P2 and P3 hold under the corrected
theorem's explicit hypotheses.

These static results do not
guarantee capacity feasibility away from equilibrium, truthful type
revelation, incentives to follow a prescribed learning algorithm, or
uniqueness of the complete quantity--price profile.

\section{The learning procedure and its limits}
Q7 proposes a variable sample-size proximal best-response algorithm with
a relaxed update. The network manager coordinates communication; each
player solves its own optimization problem using its private valuation
samples and local constraints. At iteration $i$:
\begin{enumerate}
\item Players send their current messages $s_n^i$ to the manager, who
computes and broadcasts the aggregate $\bar s^i=\sum_n s_n^i$.
Player $n$ obtains the opponents' aggregate as $\bar s^i-s_n^i$.
\item Player $n$ draws $Q_i$ independent valuation samples and computes
\[
 \widehat T_n(s^i)\in\arg\max_{z_n=(y_n,r_n)\in\mathcal S_n^P}
 \left\{\frac1{Q_i}\sum_{q=1}^{Q_i}\psi_n(y_n;\xi_{n,i}^q)
 -t_n(z_n,s_{-n}^i)-\frac\mu2\|z_n-s_n^i\|^2\right\}.
\]
Here $\mu>0$, and $\mathcal S_n^P=X_n\times[0,P_{\max}]$ includes
the coordinatewise price cap imposed in the learning section.
Opponent-only payment terms can be dropped from this optimization,
so only aggregate opponents' quantities and prices are needed.
\item All players update simultaneously, with $0<\tau<1$:
\[
 s_n^{i+1}=(1-\tau)s_n^i+\tau\widehat T_n(s^i).
\]
\end{enumerate}
The proximal penalty discourages large moves; relaxation takes only a
fraction of the computed move. Increasing $Q_i$ reduces sampling error.
The paper's convergence theorem uses geometrically increasing sample
sizes $Q_i=\lceil C_b/\eta^{2(i+1)}\rceil$, with $0<\eta<1$ and
an error-bound constant $C_b>0$. It assumes bounded gradient noise,
smoothness, a compact price domain and nonexpansiveness of the exact
proximal best-response map. Nonexpansiveness is an additional assumption,
not a consequence of the static repair theorem proved above.

Deleting the fourth payment term leaves these optimization problems
unchanged, hence leaves the intended learning updates unchanged for fixed
samples and initialization. It does not repair the convergence proof.
In particular, the paper's final convergence argument invokes uniqueness
of the Nash equilibrium, which the price-continuum example disproves.
Nonuniqueness alone does not prevent convergence to some equilibrium,
but convergence to a uniquely determined message profile is unavailable.
The algorithm also prescribes how players should update; equilibrium
incentives alone do not prove that following this entire protocol is
strategically optimal. These issues require a separate dynamic analysis.

\section{Collusion and Sybil attacks}
Nash equilibrium protects against unilateral deviations with the other
players' messages fixed. It does not protect against coordinated changes
by a coalition. Q7 uses ``strong Nash implementation'' to mean that every
Nash equilibrium implements an efficient allocation; this does not assert
that those equilibria resist coalition deviations.

\paragraph{A profitable coalition deviation in the formal game.}
Set $N=3$, $K=1$, $c=\gamma=1$, and, for every player,
\[
 X_n=[0,1],\qquad V_n(x)=2x-\tfrac12x^2.
\]
For both the original and repaired payment, the symmetric profile
\[
 x_n^*=\tfrac13,\qquad p_n^*=\tfrac53
\]
is a Nash equilibrium with zero payments. Indeed, each player's own
quantity and price gradients vanish, and its own utility Hessian is
\[
 \begin{pmatrix}-1&1/2\\1/2&-1\end{pmatrix}\prec0,
\]
so these are global best responses. Each player receives
$V_n(1/3)=11/18$.
If all players jointly change their messages to $x_n=1$, $p_n=0$,
every payment term vanishes and each player's utility becomes
$V_n(1)=3/2>11/18$. Thus this equilibrium is not resistant to coalition
deviations, for either payment rule; no side transfers are needed.

This deviation demands three units when capacity is only one. It is
admissible in the formal message game because aggregate capacity does
not constrain messages and the allocation rule is $y_n(s)=x_n$.
The example therefore exposes the importance of feasibility away from
equilibrium. It does not establish a profitable deviation in a different
game that enforces capacity for every message profile. An actual manager
would need a rationing or rejection rule, and the resulting game would
require a new equilibrium and coalition analysis.

\paragraph{Sybil resistance is not established.}
The model fixes $N$ and the players' identities; it does not allow a real
participant to create additional identities. Such identities would alter
$N$, the reference shares $c/N$, and the price aggregates in the payment
rule. To analyze this attack, one must specify how the real participant's
valuation and consumption constraints apply across identities and how
its payments are combined. Resistance to profitable use of fictitious
identities is usually called \emph{false-name-proofness}.
Neither Q7's Nash guarantees nor the static repair theorem establish this
property. This is an unmodeled attack, rather than a demonstrated
profitable Sybil deviation. Deleting the fourth term supplies no identity
control and does not resolve either the coalition counterexample or the
missing Sybil analysis.

\begin{thebibliography}{1}
\bibitem{q7}
S. Garjani, M. Shokri and H. Kebriaei.
\emph{A Systematic Approach to Mechanism Design with Stochastic Dynamic Stability}.
arXiv:2608.29130, 2026.
\url{https://arxiv.org/abs/2608.29130}.
\end{thebibliography}
\end{document}
